
	
	
	\section{静电场的高斯定理}
	
	\begin{definition}[高斯定理]
		在静电场中，通过任意闭合曲面的电场强度通量等于该曲面内电荷量的代数和除以$\varepsilon _{0}$
		
		$$\oint _{S}E\cdot dS=\frac{1}{\varepsilon _{0}}\sum _{i}q_{i}$$
		
	\end{definition}
	
	\begin{note}
		
		闭合球面上的$E$通量与球面所围成的电荷量成正比，与半径无关。
		
		高斯定理说明了电场强度对闭合面通量和场源电荷关系。
		
		闭合曲面外电荷对闭合曲线上各个点电场强度有作用，但对闭合曲面整体$E$通量贡献为$0$。
		
		适用范围：静电场、运动电荷、迅速变化的电场。
	\end{note}
	
应用方法：

一、电荷分布具有某些特殊的对称性，使相应的电场分布具有一定对称性：
\begin{example}
	1:球对称性：点电荷、均匀带点球体。2：轴对称性：均匀带电直线、柱。3：平面对称性：无限大均匀带点平板。
\end{example}

二、高斯面选择要求：

	电场强度垂直于闭合面且大小处处相等/一部分电场强度垂直于闭合面且处处相等，另一部分电场强度与闭合面平行。

\begin{example}
	球对称电场：球半径为R
	
1、$q$均匀分布在整个球体
	
（1）$P$点在球内
	
$$E = \frac{q}{4\pi\varepsilon _{0}r^2}e_r$$

（2）$P$点在球外

$$E = \frac{qr}{4\pi\varepsilon _{0}R^3}e_r$$
2、$q$均匀分布在球面上
（1）$P$点在球内

$$E =0$$

（2）$P$点在球外

$$E = \frac{qr}{4\pi\varepsilon _{0}R^3}e_r$$


\end{example}

\begin{example}
	柱对称电场：
	
（1）$P$点在圆柱内

$$ E=\frac{\lambda}{2\pi \varepsilon _{0}r}e_{r}$$

（2）$P$点在圆柱外

$$E =\frac{\lambda r}{2\pi \varepsilon _{0}R^2}e_{r}$$
\end{example}

\begin{example}
	平面对称电场：
	
$$E = \frac{\sigma}{2\varepsilon_0} e_n$$

一对电荷面密度相等异号的均匀带电无限大电场（电容板）电场强度大小为

$$E = \frac{\sigma}{\varepsilon_0} $$
\end{example}
